\section{Lead-specific missingness and asymmetric easy risk}
The learned utility combines a candidate's predicted observed gain
with its exact missing-step geometric debit. We predict annotation
availability for each of the 12 future steps from history. Replacing
these probabilities with their per-row mean is a flat-visibility
control that preserves the mean visible count. A second ablation
triples the Huber weight on underestimates of observed easy excess.
The risk head shares its initialization and draws with the earlier
conditional head. The three new heads use \LeadRiskAblUpdates{} updates
across the source rotations and seeds.

The Brier reduction from lead-specific prediction is \LeadRiskAblBrier{}.
At \LeadRiskAblEdits{} exact matched edits, the direct conservative
utility contrast is \LeadRiskAblLead{} points with ordinary conditional
risk and \LeadRiskAblTailLead{} with asymmetric risk. The asymmetric
risk contrast at fixed lead-specific gain is \LeadRiskAblTail{}.
Intervals use paired source bootstrap draws after averaging seeds.
The asymmetric head increases recording passes from \LeadRiskAblBasePass{}
to \LeadRiskAblTailPass{}/\LeadRiskAblViews{} and reduces known easy-risk
violations from \LeadRiskAblBaseEasy{} to \LeadRiskAblTailEasy{}, but
worst signed known easy degradation increases from
\LeadRiskAblBaseWorst{}\% to \LeadRiskAblTailWorst{}\%.
Calibration qualifies \LeadRiskAblQual{}/9 cells.

Of the paired missing-step debit when using lead rather than flat
probabilities, \LeadRiskAblLateShare{}\% is at steps 9--12. We therefore
compare a fixed late taper with uniform scaling on the same recipients,
matching the sum of displacements from CV for each edit. It is a
secondary diagnostic on the already exposed fitting sources.
Taper-minus-uniform conservative contrasts are \LeadRiskAblTaperCond{}
and \LeadRiskAblTaperTail{} points under the two risk heads. Matching
the sum of displacements does not match their peak or qualify a policy.
